\section*{Overview}

The Z-function stores, for every position \(i\), the length of the longest prefix of the string that matches the
suffix starting at \(i\). It is one of the cleanest linear-time string tools because its meaning is direct: ``how much
does the whole string match here?''

\section*{When to Use It}

Use the Z-function when:

\begin{itemize}
  \item you need pattern matching on \(p + \# + t\),
  \item the problem asks about borders, repetitions, or periodicity,
  \item you want a prefix-matching tool that is often easier to reason about than KMP for some tasks.
\end{itemize}

\section*{Core Idea}

Let \(z[i]\) be the longest common prefix of \(s\) and \(s[i \dots]\). The naive computation is quadratic, but the
linear algorithm keeps the rightmost matched segment \([l, r]\) and reuses previous information inside that window.

\section*{Key Insight}

Inside the current matched segment, the string already mirrors the prefix. That lets me copy part of a previous answer
instead of comparing from scratch, and only extend when necessary.

\section*{Operations / Main Technique}

\begin{itemize}
  \item compute the Z-array in one scan,
  \item pattern matching by running Z on \(p + \# + t\),
  \item detect periods or borders from large Z-values.
\end{itemize}

\paragraph{Worked example.}
For \(\texttt{s = "aaaaa"}\), the Z-array is \([0, 4, 3, 2, 1]\). Every suffix is a prefix again, just shorter.

\section*{Correctness Intuition}

The maintained window \([l, r]\) is a segment known to match the prefix. If \(i\) lies inside it, the corresponding
prefix position already tells me a lower bound for \(z[i]\). Any extra matching beyond that lower bound is found by
explicit extension. No character comparison is wasted too many times, which gives linear time.

\section*{Complexity Analysis}

The Z-function is computed in \(O(n)\) time and \(O(n)\) memory.

\section*{Implementation}

The code provides the standard linear routine and one helper that returns all pattern-occurrence positions inside a text
using the concatenation trick.

\section*{Common Pitfalls}

\begin{itemize}
  \item Forgetting to use a separator that cannot appear in the original strings.
  \item Misreading \(z[i]\) as a substring-substring LCP instead of prefix-suffix LCP.
  \item Confusing the Z-function with the prefix function from KMP.
\end{itemize}

\section*{Variants / Extensions}

\begin{itemize}
  \item Border and periodicity checks.
  \item String compression problems using smallest period detection.
  \item Comparison with KMP: same asymptotic power, different viewpoint.
\end{itemize}

\section*{Practice Problems}

\begin{itemize}
  \item Find all occurrences of a pattern in a text.
  \item Determine the smallest period of a string.
  \item Count borders or repeated prefix occurrences.
\end{itemize}

\section*{References}

\begin{itemize}
  \item \href{https://cp-algorithms.com/string/z-function.html}{cp-algorithms: Z-Function}.
  \item \href{https://usaco.guide/adv/string-search?lang=cpp}{USACO Guide: String Search}.
  \item \href{https://codeforces.com/catalog}{Codeforces Catalog}.
\end{itemize}
